% RESUMEN

\begin{abstract}

Este informe presenta la solución de la evaluación TEST2, compuesta por
tres puntos. En el primero se analiza una señal de modulación de amplitud
(AM) almacenada en un archivo CSV: a partir de su espectro se midieron una
portadora de 20~kHz y un tono transmitido de 3~kHz (con bandas laterales
en 17 y 23~kHz), y se recuperó la envolvente mediante la magnitud de la
transformada de Hilbert seguida de filtros FIR e IIR pasa-bajas. Se
explica por qué un filtro centrado en la frecuencia del mensaje aplicado
directamente sobre la señal modulada no recupera la senoidal. En el
segundo punto se decodifica un mensaje transmitido por modulación de
encendido y apagado (OOK) a 1~kbit/s sobre una portadora de 10~kHz: la
envolvente de Hilbert se suaviza, se promedia por bit y se umbraliza,
obteniendo 144 bits que forman 18 bytes ASCII y la frase
``El Rey Leon (1994)''. En el tercer punto se ejecuta un código de filtrado
espacial con núcleos de $5\times5$ sobre una fotografía real, se justifica
que la convolución 2D es un filtro y se estudia el efecto de un núcleo
cuadrado de suma cero. Todos los valores y figuras provienen de la
ejecución real del código sobre los archivos entregados.

\end{abstract}

\begin{IEEEkeywords}

Modulación AM, modulación OOK, transformada de Hilbert, filtros FIR,
filtros IIR, decodificación ASCII, convolución 2D, filtrado espacial.

\end{IEEEkeywords}
